% !TEX TS-program = pdflatex
% !TEX encoding = UTF-8 Unicode
% !BIB TS-program = bibtex
%
% ============================================================================
% paper_si.tex -- Supplementary Information for "A universal spectral lower
% bound on the leave-node-out generalisation error of graph-supervised learning"
%
% Compiles independently with article class + references/references.bib.
% ============================================================================
\documentclass[11pt,a4paper]{article}

\usepackage[margin=2.5cm]{geometry}
\usepackage[T1]{fontenc}
\usepackage[utf8]{inputenc}
\usepackage{microtype}
\usepackage{amssymb,amsmath,amstext,amsthm,mathtools}
\usepackage{bm}
\usepackage{graphicx}
\usepackage{booktabs,array,float}
\usepackage{enumitem}
\setlist{topsep=2pt,itemsep=2pt,parsep=0pt}
\usepackage{hyperref}
\usepackage{xcolor}
\hypersetup{colorlinks=true,linkcolor=blue!50!black,citecolor=blue!50!black,urlcolor=blue!50!black}

\usepackage[authoryear,round]{natbib}
\bibliographystyle{plainnat}

\graphicspath{{figures/}}

% SI appendix-style numbering : Appendix A, A.1, etc.
\renewcommand{\thesection}{\Alph{section}}
\renewcommand{\thesubsection}{\Alph{section}.\arabic{subsection}}
\renewcommand{\theequation}{\Alph{section}.\arabic{equation}}
\renewcommand{\thetable}{\Alph{section}.\arabic{table}}
\renewcommand{\thefigure}{\Alph{section}.\arabic{figure}}

\theoremstyle{plain}
\newtheorem{theorem}{Theorem}
\newtheorem{proposition}[theorem]{Proposition}
\newtheorem{lemma}[theorem]{Lemma}
\newtheorem{corollary}[theorem]{Corollary}
\theoremstyle{definition}
\newtheorem{definition}[theorem]{Definition}
\newtheorem{assumption}[theorem]{Assumption}
\theoremstyle{remark}
\newtheorem{remark}[theorem]{Remark}
\newtheorem{conjecture}[theorem]{Conjecture}
\renewcommand{\thetheorem}{\Alph{section}.\arabic{theorem}}

% Notation -- match paper.tex
\newcommand{\bphi}{\boldsymbol{\phi}}
\newcommand{\by}{\mathbf{y}}
\newcommand{\bP}{\mathbf{P}}
\newcommand{\bB}{\mathbf{B}}
\newcommand{\bU}{\mathbf{U}}
\newcommand{\bI}{\mathbf{I}}
\newcommand{\bM}{\mathbf{M}}
\newcommand{\bv}{\mathbf{v}}
\newcommand{\bu}{\mathbf{u}}
\newcommand{\bc}{\mathbf{c}}
\newcommand{\bX}{\mathbf{X}}
\newcommand{\bD}{\mathbf{D}}
\newcommand{\bW}{\mathbf{W}}
\newcommand{\bone}{\mathbf{1}}
\newcommand{\E}{\mathbb{E}}
\newcommand{\R}{\mathbb{R}}
\newcommand{\PR}{\mathbb{P}}
\newcommand{\Var}{\mathrm{Var}}
\newcommand{\Cov}{\mathrm{Cov}}
\newcommand{\tr}{\mathrm{tr}}
\newcommand{\Rspec}{R^2_{\mathrm{spec}}}
\newcommand{\Lsym}{\mathcal{L}_{\mathrm{sym}}}
\newcommand{\Scal}{\mathcal{S}}
\newcommand{\Acal}{\mathcal{A}}
\newcommand{\Hcal}{\mathcal{H}}
\newcommand{\Tcal}{\mathcal{T}}
\newcommand{\Dcal}{\mathcal{D}}
\newcommand{\Ycal}{\mathcal{Y}}
\newcommand{\Rcal}{\mathcal{R}}
\newcommand{\frR}{\mathfrak{R}}

% Cross-refs to the main paper -- rendered as visible markers.
\newcommand{\maintext}[1]{(main text, Sec.~#1)}

\title{Supplementary Information for\\
``A universal spectral lower bound on the leave-node-out generalisation
error of graph-supervised learning''}

\author{Rohan Foss\'e\thanks{rfosse@cesi.fr; ORCID 0009-0002-2195-0198.}
\and Ga\"el Pallares\thanks{ORCID 0009-0002-8680-604X.}\\
\small CESI LINEACT (EA 7527), Montpellier, France}

\date{Companion to paper.tex working draft v0.4 --- May 2026}

\begin{document}

\maketitle

\begin{abstract}
This supplementary information accompanies the manuscript
\emph{A universal spectral lower bound on the leave-node-out
generalisation error of graph-supervised learning} \maintext{1--8}.
It provides four appendices: (A) the arbitrary-subspace Pesenson
extension lemma used as a black-box reference in the proof of
Theorem~2 of the main paper; (B) the transductive Rademacher
constants of \citet{ElYanivPechyony2006} with detailed derivation;
(C) full proofs of the negative-transfer (C2) and active-learning
(C3) corollaries, with the sampling-restriction correction made
explicit; and (D) the empirical reproducibility map across the
sibling \texttt{cycling-data-lab} repositories.  An appendix~E
compares our finite-population empirical-risk setting to the
continuous-Hilbert-space noisy inverse-problem setting of
\citet{BauerPereverzev2005}, clarifying why their minimax theory
does not directly apply to our problem.
\end{abstract}

% ============================================================================
\section{Arbitrary-subspace Pesenson reconstruction lemma}
\label{sec:si-A}

The proof of Theorem~2 \maintext{5} cites the
sampling-quality-regularised reconstruction lemma of
\citet{Pesenson2008}/\citet{AnisGaddeOrtega2016} as a black-box, with
the extension from bandlimited subspaces $V_K := \mathrm{span}(\bu_1,
\ldots, \bu_K)$ to arbitrary feature subspaces $\Scal = \mathrm{col}(\bB)$
attributed to \citet{ChepuriLeus2018} and \citet{TanakaEldar2020}.
For self-containment, we record the statement and a 3-line proof.

\begin{lemma}[Subspace Pesenson reconstruction]
\label{lem:si-pesenson-subspace}
Let $\bB \in \R^{N \times d}$ be any orthonormal basis matrix
($\bB^\top \bB = \bI_d$) for a closed linear subspace $\Scal \subseteq
\R^N$.  Let $T \subset V$ with $\sigma_T(\Scal) := \lambda_{\min}(\bB^\top
\bP_T \bB) > 0$, where $\bP_T$ is the diagonal projector onto $T$.
For any $\by \in \R^N$ and any $\tau > 0$, the Tikhonov-regularised
reconstruction
\begin{equation}
    \hat\by_\tau \;:=\; \bB \big(\bB^\top \bP_T \bB + \tau \bI_d\big)^{-1}
        \bB^\top \bP_T \by
    \label{eq:si-pesenson-ridge}
\end{equation}
satisfies
\begin{equation}
    \big\|\hat\by_\tau - \bP_\Scal \by\big\|^2
    \;\leq\; \frac{\tau^2}{(\sigma_T(\Scal) + \tau)^2}\,\|\bP_\Scal\by\|^2
        \;+\; \frac{1}{(\sigma_T(\Scal) + \tau)^2}\,\|\bP_T \bP_{\Scal^\perp}\by\|^2.
    \label{eq:si-pesenson-bound}
\end{equation}
The bound does not depend on $\bB$ being a basis of Laplacian
eigenvectors; the classical \citet{Pesenson2008} statement is the
specialisation $\bB = \bU_K$.
\end{lemma}

\begin{proof}
Let $\bM := \bB^\top \bP_T \bB \in \R^{d \times d}$, symmetric PSD with
$\lambda_{\min}(\bM) = \sigma_T(\Scal) =: \sigma$.  Let $\bc^\star :=
\bB^\top \by$ so that $\bP_\Scal\by = \bB \bc^\star$.  Decompose
$\by = \bP_\Scal\by + \bP_{\Scal^\perp}\by$ and compute
\begin{equation*}
    \bB^\top \bP_T \by = \bB^\top \bP_T \bB \bc^\star + \bB^\top \bP_T \bP_{\Scal^\perp}\by
    = \bM \bc^\star + \bB^\top \bP_T \bP_{\Scal^\perp}\by.
\end{equation*}
Hence the Tikhonov estimate in coefficient space is
$\hat\bc_\tau = (\bM + \tau\bI)^{-1}\bB^\top \bP_T \by$, and
\begin{align*}
    \hat\bc_\tau - \bc^\star
    &= \big[(\bM + \tau\bI)^{-1}\bM - \bI\big]\bc^\star
        + (\bM + \tau\bI)^{-1}\bB^\top \bP_T \bP_{\Scal^\perp}\by\\
    &= -\tau (\bM + \tau\bI)^{-1}\bc^\star
        + (\bM + \tau\bI)^{-1}\bB^\top \bP_T \bP_{\Scal^\perp}\by.
\end{align*}
Using $\|(\bM + \tau\bI)^{-1}\|_{\mathrm{op}} \leq 1/(\sigma + \tau)$
and $\|\bB^\top\|_{\mathrm{op}} = 1$,
\begin{equation*}
    \|\hat\bc_\tau - \bc^\star\|^2 \leq \frac{\tau^2}{(\sigma + \tau)^2}\|\bc^\star\|^2 + \frac{1}{(\sigma + \tau)^2}\|\bP_T \bP_{\Scal^\perp}\by\|^2.
\end{equation*}
Since $\bB$ is an isometry, $\|\hat\by_\tau - \bP_\Scal\by\| = \|\hat\bc_\tau - \bc^\star\|$ and $\|\bc^\star\| = \|\bP_\Scal\by\|$, yielding \eqref{eq:si-pesenson-bound}.\qed
\end{proof}

\begin{remark}[Where Lemma~\ref{lem:si-pesenson-subspace} would enter the
main-paper proof]
\label{rem:si-when-pesenson}
The main-paper proof of Theorem~2 does \emph{not} actually invoke
Lemma~\ref{lem:si-pesenson-subspace}, because the witness predictor
used there is $f^\star_{\Hcal_d}$ (the population-risk minimiser), which
requires no Tikhonov regularisation.  The lemma is preserved here as a
reference because: (i) earlier drafts of the program
(\texttt{notes/01\_universal\_bound\_proof.tex}, since superseded)
routed through Pesenson-ridge as the witness; (ii) it is the natural
substitute when $\bP_\Scal\by \notin \Hcal_d$ (non-realisable case),
treated in \texttt{notes/03\_non\_realisable\_resolution.tex}; (iii)
it is the foundation of the leverage-score active-learning analysis
(C3).
\end{remark}

% ============================================================================
\section{Transductive Rademacher constants}
\label{sec:si-B}

The slack of Theorem~2 of the main paper relies on the transductive
Rademacher bound of \citet[][Theorem~1]{ElYanivPechyony2006}, which we
restate here in our notation.

\begin{theorem}[Transductive Rademacher bound]
\label{thm:si-trans-rademacher}
Let $N$, $n$, $u = N - n$ as in \maintext{3}.  Let $\ell : \R \to [0,
B]$ be a loss function and $\Hcal_d \subseteq \R^N$ a hypothesis
class.  Define the transductive Rademacher complexity
\begin{equation}
    \mathcal{R}_n^{\mathrm{trans}}(\Hcal_d)
    \;:=\; \frac{1}{n}\Big(\frac{1}{n} + \frac{1}{u}\Big)^{-1/2}
        \;\E_\sigma\!\!\left[\sup_{f \in \Hcal_d}\,\sum_{v \in V}\sigma_v f(v)\right],
    \label{eq:si-trans-rademacher-defn}
\end{equation}
where $\sigma_v$ are i.i.d.\ taking $+1$ w.p.\ $p = n/N$, $-1$ w.p.\
$p$, and $0$ w.p.\ $1 - 2p$.  Then for any learner $\hat f$ outputting
predictors in $\Hcal_d$, with probability $\geq 1 - \eta$ over the
uniformly random split $(T, T^c)$ of $V$ at $|T| = n$,
\begin{equation}
    L_{T^c}(\hat f) \;\leq\; L_T^\pi(\hat f)
        + 2 \mathcal{R}_n^{\mathrm{trans}}(\ell \circ \Hcal_d)
        + B \sqrt{\frac{(n + u)\log(1/\eta)}{n u}}
        \cdot \big(11/4 + \sqrt{\log(1/\eta)/2}\big),
    \label{eq:si-trans-rademacher-bound}
\end{equation}
where $L_T^\pi(\hat f) := \frac{1}{N}\sum_{v \in T}\ell(\hat f(v) - y_v)/\pi(v)$
is the importance-weighted training loss (which under uniform
sampling reduces to $L_T(\hat f) \cdot N/n$).
\end{theorem}

The constant $11/4 + \sqrt{\log(1/\eta)/2}$ is bounded by $\approx 5.05$
for $\eta = 0.05$, which is the value used in the main paper.

The key structural property is that the bound \emph{directly}
controls $L_{T^c} - L_T^\pi$ without an intermediate population-risk
term $L_V$.  Routing through $L_V$ would incur a factor of $\rho(\Pi)
= N/(N-n)$ on the slack (transduction identity), inflating the slack
by $\sqrt{\rho(\Pi) - 1}$ relative to
\eqref{eq:si-trans-rademacher-bound}; see
\texttt{notes/04\_sharp\_transductive\_constant.tex}~\S 1--3 for the
detailed comparison.

\begin{lemma}[Talagrand contraction for the transductive Rademacher]
\label{lem:si-talagrand-trans}
For the $2M$-Lipschitz square loss truncated to $[-M, M]$ and any
hypothesis class $\Hcal_d \subseteq [-M, M]^N$,
\begin{equation*}
    \mathcal{R}_n^{\mathrm{trans}}(\ell \circ \Hcal_d) \;\leq\; 4 M \cdot \mathcal{R}_n^{\mathrm{trans}}(\Hcal_d).
\end{equation*}
\end{lemma}

\begin{proof}
\citet[][\S 4.2]{Talagrand1996concentration} for the standard Rademacher,
extended to the transductive case by an analogous symmetrisation
argument
(\citealp[][\S 5]{ElYanivPechyony2006}; the only change is that the
$\sigma_v$ take value $0$ on the $1 - 2p$ portion of $V$).
\end{proof}

% ============================================================================
\section{Negative-transfer corollary: technical conditions}
\label{sec:si-C}

We make explicit the sampling-restriction correction of the negative-transfer
corollary (C2 in the main paper, Remark on
``sampling-restriction correction'').

\subsection{Setup}

Let $V = V_{\mathrm{src}} \cup V_{\mathrm{tgt}}$ disjoint, $|V_{\mathrm{src}}| =
n_s$, $|V_{\mathrm{tgt}}| = n_t$, $N = n_s + n_t$.  Let $\Hcal_d$ be a
closed linear subspace of $\R^N$ of dimension $d$, with orthonormal
basis $\bB \in \R^{N \times d}$.

Define the source-set and target-set sampling-quality scalars
\begin{align}
    \sigma_{V_{\mathrm{src}}}(\Hcal_d) &:= \lambda_{\min}\big(\bB^\top \bP_{V_{\mathrm{src}}}\bB\big), &
    \sigma_{V_{\mathrm{tgt}}}(\Hcal_d) &:= \lambda_{\min}\big(\bB^\top \bP_{V_{\mathrm{tgt}}}\bB\big).
\end{align}

\subsection{Identifiability of the source ERM}

\begin{lemma}[Source identifiability]
\label{lem:si-source-identif}
If $\sigma_{V_{\mathrm{src}}}(\Hcal_d) > 0$ (i.e., $V_{\mathrm{src}}$ is a
\emph{uniqueness set} for $\Hcal_d$), then the source-ERM
\[
    \hat\bc_s := \arg\min_{\bc \in \R^d} \frac{1}{n_s}\sum_{v \in V_{\mathrm{src}}}\big(\langle\bB(v, :), \bc\rangle - y_v\big)^2
\]
is unique and $\hat\bc_s = (\bB^\top \bP_{V_{\mathrm{src}}}\bB)^{-1}\bB^\top \bP_{V_{\mathrm{src}}}\by$.
\end{lemma}

\begin{proof}
Standard OLS on the design matrix $\bB[V_{\mathrm{src}}, :] \in
\R^{n_s \times d}$, which has rank $d$ iff
$\sigma_{V_{\mathrm{src}}}(\Hcal_d) > 0$.
\end{proof}

\subsection{Inner-product approximation on $V_{\mathrm{tgt}}$}

\begin{lemma}[Restriction-corrected orthogonality]
\label{lem:si-restriction-orth}
Let $f, g \in \Hcal_d$.  Then
\begin{equation}
    \Big|\langle f, g\rangle_{V_{\mathrm{tgt}}} - \tfrac{n_t}{N}\langle f, g\rangle\Big| \;\leq\; M^2 \cdot \big(1 - \sigma_{V_{\mathrm{tgt}}}(\Hcal_d) \tfrac{N}{n_t}\big)_+,
\end{equation}
where the subscript-$+$ denotes positive part.  In particular, if
$V_{\mathrm{tgt}}$ is a uniqueness set ($\sigma_{V_{\mathrm{tgt}}} \cdot N/n_t \geq 1$),
the correction vanishes and $\langle f, g\rangle_{V_{\mathrm{tgt}}} = (n_t/N)\langle f, g\rangle$.
\end{lemma}

This lemma is what makes the negative-transfer proof (main paper,
Theorem~6) clean: under the assumption that $V_{\mathrm{tgt}}$ is a
uniqueness set for $\Hcal_d$ (which is implicit in evaluating any
non-degenerate target loss), the orthogonality of source and target
projections under spectral disjointness translates from the global
$\R^N$ inner product to the restricted $V_{\mathrm{tgt}}$ inner product.

\begin{remark}[When does $V_{\mathrm{tgt}}$ fail to be a uniqueness set?]
\label{rem:si-uniqueness-failure}
$V_{\mathrm{tgt}}$ fails to be a uniqueness set for $\Hcal_d$ iff
$\dim(\Hcal_d) > n_t$, i.e., the target sample is smaller than the
hypothesis class dimension.  In this regime the target ERM is not
identifiable, and the negative-transfer phenomenon is masked by the
broader unidentifiability of any prediction on $V_{\mathrm{tgt}}$
(regardless of source). We assume $n_t \geq d$ throughout the
corollary; in practice this is mild (it requires more target
examples than features), and when it fails the more fundamental
identifiability question must be addressed first.
\end{remark}

% ============================================================================
\section{Reproducibility map}
\label{sec:si-D}

The empirical content of the main paper
(\maintext{8}, Table~1) is reproducible from publicly archived
scripts in the sibling \texttt{cycling-data-lab} repositories.  We
provide a per-row mapping.

\begin{table}[H]
\centering
\small
\begin{tabular}{lll}
\toprule
\textbf{Main-paper Table~1 row} & \textbf{Repository} & \textbf{Script (\texttt{experiments/})} \\
\midrule
8 MatBench tasks (mp\_e\_form, mp\_gap, etc.) & \texttt{materials-applicability-bound} & \texttt{d11\_predictive\_multitask.py} \\
matbench\_expt\_gap                            & \texttt{materials-applicability-bound} & \texttt{d11b\_expt\_gap.py} \\
MovieLens user-rating pilot                    & \texttt{materials-applicability-bound} & \texttt{d01\_movielens\_pilot.py} \\
QSAR pilot                                     & \texttt{materials-applicability-bound} & \texttt{d04\_qsar\_pilot.py} \\
AFLOW bulk / shear modulus pilots              & \texttt{materials-applicability-bound} & \texttt{d03\_aflow\_pilot.py} \\
Bike-share LSO panel ($\Delta R^2 \in [0.06, 0.55]$) & \texttt{bikeshare-demand-forecasting} & \texttt{d19\_lso\_rigorous.py} \\
Mobility commune panel (in preparation)        & \texttt{mobility-applicability-bound} & (in preparation) \\
\midrule
Falsifiability F1 (Spearman exact $p = 4.96 \times 10^{-5}$) & \texttt{materials-applicability-bound} & \texttt{d20\_exact\_permutation\_test.py} \\
Falsifiability F2 (partial Spearman, bootstrap CI)           & \texttt{materials-applicability-bound} & \texttt{d13\_bootstrap\_ci.py} \\
Falsifiability F3 (CIG, shuffled-$k$NN null)                 & \texttt{materials-applicability-bound} & \texttt{d25\_knn\_sensitivity.py}, \texttt{d25b\_*} \\
Spectrum decomposition (Magpie / chem graph)                 & \texttt{materials-applicability-bound} & \texttt{d05\_rspec\_matbench.py}, \texttt{d06\_h1\_spectrum\_matbench.py} \\
CHGNet encoder discrimination                                & \texttt{materials-applicability-bound} & \texttt{d24\_chgnet\_structural\_eg.py} \\
Figure~1 (spectral-projection picture)                       & \texttt{structural-bounds-framework} & \texttt{d01\_spectral\_projection\_figure.py} \\
\bottomrule
\end{tabular}
\caption{Per-row reproducibility map for the main-paper empirical
section.  Every value is regenerable from raw open data via the
indicated script, using SEED = 42 and Python 3.12 with the
\texttt{requirements.txt} of the corresponding repository.}
\label{tab:si-reproducibility}
\end{table}

\subsection{Random seeds and finite-population corrections}

The master random seed is pinned to \texttt{SEED = 42} in every
stochastic script.  Per-seed sensitivity is documented in
\texttt{d18\_shuffled\_knn\_null.py} (10 realisations) and
\texttt{d25\_knn\_sensitivity.py} (8 tasks $\times$ 3 $k$ values).
The headline Spearman $\rho = +1.000$ is robust to the choice of
$k$ across $k \in \{5, 10, 20\}$ (see \texttt{outputs/d25\_knn\_sensitivity\_full.json}).

Finite-population corrections in the without-replacement Hoeffding
inequality (Serfling 1974) are made explicit in the proofs of
Theorems~1 and~2 of the main paper; the corresponding empirical
correction factor $1 - (n-1)/N$ is applied in
\texttt{d13\_bootstrap\_ci.py}.

% ============================================================================
\section{Comparison to noisy continuous inverse-problem theory}
\label{sec:si-E}

A natural question is why our framework does not invoke the classical
minimax theory of noisy linear inverse problems on continuous Hilbert
spaces \citep{BauerPereverzev2005,Cavalier2008}.  We summarise the
mismatch.

\paragraph{Setting comparison.}

\begin{center}
\begin{tabular}{lll}
\toprule
\textbf{Feature}         & \textbf{Bauer--Pereverzev 2005}            & \textbf{Present paper} \\
\midrule
Domain                   & continuous Hilbert space                    & finite graph $V$, $|V| = N$ \\
Forward operator         & compact linear operator $A : X \to Y$       & sample restriction $\bP_T : \R^N \to \R^n$ \\
Noise                    & Gaussian $\sigma \xi$ on observations       & none (deterministic target) \\
Source of randomness     & noise $\xi$                                 & random sub-sample $T$ \\
Estimand                 & full signal $f \in X$                       & held-out loss $L_{T^c}(\hat f)$ \\
Minimax-optimal estimator & Tikhonov, $\tau \sim \sigma^{2/(2\mu+1)}$ & ERM with witness $f^\star_{\Hcal_d}$ \\
Rate                     & $\sigma^{4\mu/(2\mu+1)}$                    & $M^2/\sqrt{N-n}$ (Berry--Esseen) \\
\bottomrule
\end{tabular}
\end{center}

\paragraph{Why the rates differ.}  In \citet{BauerPereverzev2005}, the
rate $\sigma^{4\mu/(2\mu+1)}$ is the variance-bias trade-off of a
Tikhonov regulariser against Gaussian noise of variance $\sigma^2$ on
observations.  In our finite-population setting, the analogous
``noise'' would be the random restriction $\by|_T$ instead of $\by|_V$,
but this is not Gaussian: it is a uniform sub-sample.  The
Berry--Esseen rate $1/\sqrt{N-n}$ governs the residual fluctuation,
which is the right object to bound in this setting; the inverse-problem
rate $\sigma^{4\mu/(2\mu+1)}$ does not apply because there is no
parametric noise variance to optimise against.

\paragraph{When inverse-problem theory \emph{would} apply.}  If we
augment the framework with noisy observations $y_v^{\mathrm{obs}} = y_v
+ \xi_v$, $\xi_v \sim \mathcal{N}(0, \sigma_\xi^2)$ i.i.d.\ on $T$,
the Bauer--Pereverzev rate would govern the additional noise-dependent
slack term.  This is a natural extension to chemical-precision
applications (e.g., DFT energy targets with quantum-Monte-Carlo
ground-truth labels carrying explicit noise estimates) but is not
operative in the deterministic-target setting of our motivating
applications (MatBench, MovieLens, mobility panels).

% ============================================================================
\section*{References}
\bibliography{references/references}

\end{document}
